\setcounter{section}{0}

\Needspace{5\baselineskip}
\hypertarget{ux6d41ux5339ux914dux6a21ux578b-1}{%
\section{流匹配模型}\label{ux6d41ux5339ux914dux6a21ux578b-1}}

\Needspace{5\baselineskip}
\hypertarget{ux4e00ux6d41ux5339ux914dux7684ux57faux672cux601dux8defux53caux4e0eux6269ux6563ux6a21ux578bux7684ux5bf9ux6bd4}{%
\subsection{一、流匹配的基本思路及与扩散模型的对比}\label{ux4e00ux6d41ux5339ux914dux7684ux57faux672cux601dux8defux53caux4e0eux6269ux6563ux6a21ux578bux7684ux5bf9ux6bd4}}

流匹配（Flow Matching）和扩散模型都是生成模型，它们的目标都是将简单的分布（如高斯噪声，\(p_0\)）变换为复杂的数据分布（如图像或文本，\(p_1\)）。

扩散模型的思路是：像墨水在水中扩散一样，数据逐渐加噪变成噪声。生成过程则是逆转这个扩散过程，像``去噪''一样一步步还原数据。这通常被建模为随机微分方程（SDE），路径是离散、曲折的。

流匹配的思路是：希望从加噪后的分布直接得到加噪前的分布。直接定义一个向量场，这个场告诉粒子（数据点）在每一时刻 \(t\) 应该往哪个方向移动。

\Needspace{5\baselineskip}
流匹配中的几个基本概念：

\begin{itemize}
\tightlist
\item
  \textbf{起始状态}：在 \(t=0\) 时刻，所有样本都服从一个简单分布 \(p_0(x)\)，例如标准高斯噪声，这就像一团散沙。
\item
  \textbf{目标状态}：在 \(t=1\) 时刻，样本汇聚成了我们想要生成的复杂数据分布 \(p_1(x)\)，例如精美的图片集合，就像沙子堆成的城堡。
\item
  \textbf{速度场}：记为 \(v(t,x)\)，在任意时刻 \(t\in[0,1]\) 和高维空间的任意位置 \(x\)，存在一个``风向标''或者说``力场''，它指示了位于该位置的粒子在下一瞬间应该往哪个方向移动、速度有多快。
\item
\Needspace{5\baselineskip}
  \textbf{常微分方程的作用}：粒子的运动轨迹 \(x_t\) 随时间的变化率（即速度），正是由该位置的速度场决定的。这在数学上天然地被表达为一个常微分方程：
\end{itemize}

\[
\frac{dx_t}{dt}=v(t,x_t)
\]

如果我们知道每个时刻和位置的风向（即神经网络学到的速度场），我们就必须通过积分（求解这个 ODE）来推导出这粒沙子最终会被吹到哪里。

\Needspace{5\baselineskip}
我们希望让特定数据 \(a\) 从 \(x_0\)（噪声位置，已知）以给定的速度分布（也就是``速度场''中对应位置对应时间的值）趋近 \(x_1\)（有意义的位置，未知）。事实上，上述常微分方程可写成积分形式：

\[
x_1-x_0=\int_0^1 v_t\,dt
\]

由于我们不关心加噪的过程经历了什么，因此我们不妨假设一个理想情况，即所有数据都沿着我们希望的速度和方向匀速直线运动。因此，我们的目标就是要学习这个速度场，使之接近我们假设的理想情况。如果我们把大量噪声点扔进这个场里，它们最终会汇聚成有意义的数据。

\Needspace{5\baselineskip}
\hypertarget{ux4e8cux6570ux5b66ux539fux7406}{%
\subsection{二、数学原理}\label{ux4e8cux6570ux5b66ux539fux7406}}

流匹配的核心在于定义和学习这个``速度场''。

\Needspace{5\baselineskip}
\hypertarget{a.-ux6982ux7387ux8defux5f84probability-path}{%
\subsubsection{A. 概率路径（Probability Path）}\label{a.-ux6982ux7387ux8defux5f84probability-path}}

\Needspace{5\baselineskip}
我们不需要像扩散模型那样模拟复杂的物理扩散过程。我们可以直接定义一条从噪声 \(x_0\) 到数据 \(x_1\) 的路径。最简单且最高效的方法是线性插值（Linear Interpolation）：

\[
x_t=(1-t)x_0+tx_1,\quad t\in[0,1]
\]

\begin{itemize}
\tightlist
\item
  当 \(t=0\) 时，\(x_t=x_0\)，全是噪声。
\item
  当 \(t=1\) 时，\(x_t=x_1\)，全是数据。这条路径是直线的，这对于快速生成非常重要。
\end{itemize}

\Needspace{5\baselineskip}
\hypertarget{b.-ux901fux5ea6ux573avelocity-field}{%
\subsubsection{B. 速度场（Velocity Field）}\label{b.-ux901fux5ea6ux573avelocity-field}}

\Needspace{5\baselineskip}
既然有了位置 \(x_t\) 随时间 \(t\) 的变化公式，我们就可以算出粒子移动的速度（对时间求导）：

\[
u_t(x_t\mid x_0,x_1)=\frac{d}{dt}x_t=x_1-x_0
\]

这个 \(u_t\) 就是条件流匹配（Conditional Flow Matching）的目标速度。它告诉我们：如果要在时间 \(t=1\) 到达 \(x_1\)，当前的粒子应该以什么速度和方向移动。

\Needspace{5\baselineskip}
\hypertarget{c.-ux8badux7ec3ux76eeux6807flow-matching-objective}{%
\subsubsection{C. 训练目标（Flow Matching Objective）}\label{c.-ux8badux7ec3ux76eeux6807flow-matching-objective}}

在训练时，我们无法预知终点 \(x_1\)，这里只是我们想生成的。因此，我们需要训练一个神经网络 \(v_\theta(x_t,t)\)，让它在只知道当前位置 \(x_t\) 和时间 \(t\) 的情况下，去预测上述的目标速度。

\Needspace{5\baselineskip}
损失函数非常简单，就是回归（Regression）问题：

\[
\mathcal{L}_{FM}(\theta)=\mathbb{E}_{t,x_0,x_1}\left[\left\|v_\theta(x_t,t)-u_t(x_t\mid x_0,x_1)\right\|^2\right]
\]

也就是说，让神经网络预测的``流速度''尽可能接近真实的``流速度''。

当然，实际上在图像生成等任务中，预测 \(v\) 时已知的信息包括当前位置 \(x_t\)、时间 \(t\) 和条件 \(c\)。

在给定终点的情况下，我们希望速度应为指向终点的常数。或许你会问，那为什么还需要用神经网络拟合呢？这是因为我们在推理的时候并不知道终点是什么，不同的扩散起点（也就对应不同``粒子''）、不同的附加条件都会对应不同的终点和不同的``理想速度''。但在巨大的高维向量场中，实际在执行特定任务时的速度（该位置的场的方向）往往不是直指向终点的匀速，流匹配模型就是尽量去在每个给定的点拟合我们想要的速度方向。

\Needspace{5\baselineskip}
\hypertarget{ux4e09ux8badux7ec3ux65b9ux6cd5ux4ee5ux6d41ux5339ux914dvlaux6a21ux578bux4e3aux4f8b}{%
\subsection{三、训练方法（以流匹配VLA模型为例）}\label{ux4e09ux8badux7ec3ux65b9ux6cd5ux4ee5ux6d41ux5339ux914dvlaux6a21ux578bux4e3aux4f8b}}

以目前最先进的最优传输流匹配（Optimal Transport Flow Matching，类似于 \(\pi_0\) 中的设计）为例。假设基础噪声分布为 \(\pi_0\sim p_0=\mathcal{N}(0,I)\)，真实动作轨迹分布为 \(\pi_1\sim p_1\)。多模态条件为 \(c=\mathrm{Encoder}(o_t,I_t)\)。

\Needspace{5\baselineskip}
模型构建了一条连接噪声与真实动作的直线概率路径：

\[
x_t=(1-t)x_0+tx_1
\]

\Needspace{5\baselineskip}
动作生成器 \(v_\theta\) 的目标是拟合真实的速度场（Velocity Field），其目标函数定义为：

\[
\mathcal{L}_{FM}(\theta)=\mathbb{E}_{x_0,x_1,t}\left[\left\|v_\theta(x_t,t,c)-(x_1-x_0)\right\|^2\right]
\]

这里 \((x_1-x_0)\) 是真实路径的恒定速度。相较于传统的 DDPM 扩散模型，流匹配的 ODE 轨迹更加平滑且笔直，能够在较少的推理步数（甚至几步内）生成高保真的连续动作信号。

在训练时，我们对于每一组 \(x_0\) 和 \(x_1\)，都希望其连线上每个位置的速度尽可能满足从 \(x_0\) 指向 \(x_1\)，且大小在数值上等于二者之差。当然事实上，由于空间中的速度场是连续的，因此只能尽量拟合上述需求。当多条连线同时汇聚于一个点时，该点速度近似等于各连线对应速度的平均值。

\Needspace{5\baselineskip}
工作流：

\begin{enumerate}
\def\labelenumi{\arabic{enumi}.}
\tightlist
\item
  \textbf{特征提取}：VLM 接收当前相机观测 \(o_t\) 和语言 \(l\)，前向传播生成多模态条件向量 \(c\)。
\item
  \textbf{噪声初始化}：从标准正态分布中采样一段长度为 \(H\) 的轨迹高斯噪声 \(x_0\sim\mathcal{N}(0,I)\)。
\item
  \textbf{常微分方程求解（ODE Solving）}：在时间步 \(t\in[0,1]\) 之间，动作生成器 \(v_\theta(x_t,t,c)\) 结合上下文 \(c\) 计算当前状态的梯度。通过数值积分器（如 Euler 法或 Heun 法）逐步更新 \(x_t\)。
\item
  \textbf{动作执行}：将 ODE 终点 \(x_1\) 作为预测的动作块 \(A_t\)。系统采用``模型预测控制（MPC / Receding Horizon Control）''策略，仅执行前 \(k\) 个动作，随后重新获取新的视觉观测，进入下一轮闭环推理。
\end{enumerate}

这里的时间步是连续的，但神经网络只能处理离散的 \(t\) 值输入。故往往采取采样的方法，如每隔 \(0.02\) 的时间间隔进行一次采样，用该时刻网络输出的值对时间间隔积分，让 \(x_t\)``走一小步''。

\Needspace{5\baselineskip}
\hypertarget{ux56dbux5e94ux7528-1}{%
\subsection{四、应用}\label{ux56dbux5e94ux7528-1}}

\begin{enumerate}
\def\labelenumi{\arabic{enumi}.}
\item
  生成高分辨率、高质量的图像，性能可与Stable Diffusion等扩散模型媲美甚至更优，且通常生成速度更快（因为轨迹更直，从而需要的采样步数更少） 。
\item
  用于生成连续的视频帧序列，处理高维度的时空数据。
\item
  用于科学计算与模拟，如模拟物理系统（如流体动力学）、分子构象生成（蛋白质 3D 结构预测）等，因为这些问题本质上就是对连续变量变化（高维空间中点的移动速度随时空的变化）的建模。
\end{enumerate}

\FloatBarrier
\Needspace{5\baselineskip}
\hypertarget{ux53c2ux8003ux6587ux732e-135}{%
\subsection{参考文献}\label{ux53c2ux8003ux6587ux732e-135}}

\vspace{0.25em}
\begin{itemize}
\tightlist
\item
  Lipman, Y., Chen, R. T. Q., Ben-Hamu, H., Nickel, M., \& Le, M. (2023). \href{https://arxiv.org/abs/2210.02747}{Flow Matching for Generative Modeling}. ICLR.
\item
  Liu, X., Gong, C., \& Liu, Q. (2023). \href{https://arxiv.org/abs/2209.03003}{Flow Straight and Fast: Learning to Generate and Transfer Data with Rectified Flow}. ICLR.
\end{itemize}

\clearpage
